\documentclass[10pt,a4paper]{amsart}
\usepackage[margin=19mm]{geometry}
\usepackage{amsmath,amssymb,mathtools}
\usepackage[scr=rsfs,cal=euler]{mathalfa}
\usepackage{xfrac}
\usepackage[unicode,hidelinks,bookmarks=false]{hyperref}
\newcommand{\ie}{i.e.\ }
\newcommand{\cf}{cf.\ }
\newcommand{\resp}{respectively\ }
\newcommand{\Subsection}{\subsection}
\providecommand{\nobreakdash}{\nobreak\hbox{-}}
\setlength{\emergencystretch}{2em}
\multlinegap0pt
\usepackage{bm}
\usepackage{bbm}
\usepackage{dsfont}
\usepackage{xspace}
\usepackage{cancel}
\usepackage{mathtools}
\usepackage{amsthm}
\makeatletter
\@ifundefined{theorem}{\newtheorem{theorem}{Theorem}}{}
\@ifundefined{Thm}{\newtheorem{Thm}[theorem]{Theorem}}{}
\makeatother
\title[The rectangle condition does not detect the strong irreducib]{The rectangle condition does not detect the strong irreducibility}
\author{Bo-hyun Kwon, Sungmo Kang,, Jung Hoon Lee}
\date{Source: arXiv:2509.11701v1 (CC BY 4.0)}
\begin{document}
\maketitle

\noindent\emph{Source and licence.} The rectangle condition does not detect the strong irreducibility. Version arXiv:2509.11701v1 (CC BY 4.0), distributed under the Creative Commons Attribution 4.0 International licence, as recorded in the arXiv licence metadata for this exact version and shown on the abstract page https://arxiv.org/abs/2509.11701v1. Full rights notice: licenses/arxiv-2509.11701-CC-BY-4.0.txt.\par\smallskip

\noindent\textit{Editorial scope.} Counted unit: the Thm whose source label is \texttt{main result-intro-2} in the article (source0.tex lines 570--575, source label \texttt{main result-intro-2}), delivered together with the article's own complete original proof of it (source0.tex lines 578--601, one \texttt{proof} environment of 24 lines). No mathematics is written, repaired or re-derived: statement and proof are the article's own text line for line. Required context retained uncounted: none. Every definition, display and label of those lines is retained unchanged. Omitted from this package and not redistributed: 1--569, 576--577, 602--953. Nothing inside a retained excerpt is omitted or altered: every retained body is delivered exactly as the article's lines are written, so no omission receipt of any kind accompanies this package. Citations: the retained excerpts carry no citation command at all, so no bibliography entry of the article is transcribed and no reference is redistributed. Reference adaptation: none was needed and none was made; every reference inside the delivered unit resolves inside it, so no unresolved reference or citation is produced. Printed slips kept literal: no printed word, symbol, hypothesis or number of the counted statement or of its proof was altered. Layout and class: the article's own class is \texttt{amsart} with packages including geometry, latexsym, float, amssymb, xy, epsfig, color, graphics, hyperref; the wrapper is the lane's pinned \texttt{amsart} editorial wrapper at 10pt on A4, which restates the theorem-like environments the retained text needs and adds the article's own macro definitions that the retained text needs (none), with extra wrapper packages (bm, bbm, dsfont, xspace, cancel, mathtools). The delivered numbering is the unit's own. Assets: no retained span contains a figure environment, a graphics-inclusion command, a PDF-page inclusion, a table or a TikZ picture; no image file is copied into this package, the package declares no asset dependency and no image file is redistributed with it. Licence: version arXiv:2509.11701v1 of this article is distributed under the Creative Commons Attribution 4.0 International licence, as recorded in the arXiv licence metadata for this exact version and shown on the abstract page https://arxiv.org/abs/2509.11701v1. A full rights notice is delivered at licenses/arxiv-2509.11701-CC-BY-4.0.txt. Characters: every retained excerpt is pure ASCII, so the delivered engine, XeLaTeX with its default Latin Modern text font, reports no missing character.
\begin{Thm}\label{main result-intro-2}
Let $(T_1, T_2; S)$ be a 3-bridge decomposition of a knot $K$ in $S^3$ and
$(V, W;\Sigma)$ the genus 2 Heegaard splitting obtained by the double branched covering of
$(T_1, T_2; S)$. Then $(V, W;\Sigma)$ is strongly irreducible if and only if $(T_1, T_2; S)$
is strongly irreducible.
\end{Thm}

\begin{proof}
First, we prove the `only if' part, which is easier to prove. Suppose for contradiction that $(T_1, T_2; S)$
is not strongly irreducible. Then there exists a compression $D$ in $T_1$ and a compression $E$ in $T_2$ which
are disjoint. By the property of the double branched covering, each compression of $D$ and $E$ lifts to two copies of a nonseparating essential disk in $V$ and $W$, respectively.
Therefore, the corresponding essential disks in $V$ and $W$ are disjoint, contradicting the assumption that $(V, W; \Sigma)$ is strongly irreducible. Hence, $(T_1, T_2; S)$ must be strongly irreducible.

Next, we prove the `if' part. Suppose for contradiction that $(V, W;\Sigma)$ is not strongly irreducible, i.e.,
it is weakly reducible. It is well known that a genus 2 weakly
reducible Heegaard splitting is a reducible Heegaard splitting.
Therefore, $(V, W;\Sigma)$ is one of
the following:


\begin{itemize}
\item[(i)] the connected sum of two genus 1 Heegaard splittings of $S^3$,
\item[(ii)] the connected sum of a genus 1 Heegaard splitting of a lens space and a genus 1
Heegaard splitting of $S^3$,
\item[(iii)] the connected sum of two genus 1 Heegaard splittings of lens spaces.
\end{itemize}

For (i), (ii), (iii),  the corresponding $3$-bridge decomposition $(T_1, T_2; S)$ would be a $3$-bridge decomposition of an unknot, a 2-bridge
knot, a connected sum of two 2-bridge knots, respectively. For all such knots, a 3-bridge
decomposition is always weakly reducible, a contradiction.
\end{proof}

\end{document}
